\section{Related Work}
\label{sec:related}

\para{Why safety training fails.} \citet{wei2023jailbroken} attribute jailbreaks to two failure modes: competing objectives, where helpfulness and safety pull apart, and mismatched generalisation, where safety training does not cover inputs that the model's capabilities still cover. \citet{qi2025shallow} argue that alignment is often ``a few tokens deep'', concentrated on the first tokens of the response. \citet{arditi2024refusal} find that refusal is mediated by a single direction in activation space. The gradient-masking analogy we test comes from adversarial robustness, where defences that obscure gradients give a false sense of security \cite{papernot2017practical,athalye2018obfuscated}. Building on this framing, we ask a behavioural question that these works leave open: which surface features does refusal actually depend on, and does the answer change with model size?

\para{Reformulation and optimisation attacks.} Gradient-based suffixes \cite{zou2023universal}, iterative LLM attackers such as PAIR \cite{chao2023pair}, persuasive paraphrases \cite{zeng2024johnny}, and simple past-tense rewrites \cite{andriushchenko2024pasttense} all show that re-wording a harmful request can defeat refusal. These attacks change many properties of a request at once. Unlike them, we use a single neutral rewrite and pair it with a control rewrite that keeps the trigger words, so that the lexical component of the effect can be measured on its own.

\para{Multilingual jailbreaks.} \citet{yong2023lowresource} and \citet{deng2024multilingual} report much higher unsafe-response rates for low-resource languages. \citet{deng2024multilingual} also note that weaker models often produce invalid, off-topic output in these languages. We extend this observation to a size ladder of current open models and show that, for small models, non-refusal and harmful compliance differ by a factor of three to ten.

\para{Role-play and multi-turn attacks.} In-the-wild ``Do Anything Now'' prompts use personas that claim to have no rules \cite{shen2024dan}, and Crescendo builds up to a harmful request over several benign turns \cite{russinovich2024crescendo}. We use short, fixed templates of both kinds so that the same framing is applied to every model and every request.

\para{Exaggerated safety and evaluation.} \xstest \cite{rottger2024xstest} measures over-refusal of safe prompts that contain trigger words, such as ``kill a Python process''. We reuse it as the ``keywords without harm'' arm of our decorrelated pool. For evaluation, we follow \harmbench \cite{mazeika2024harmbench} and \jbb \cite{chao2024jailbreakbench}, which provide behaviour sets and a fine-tuned harmfulness classifier, and use an LLM judge in the style of \citet{zheng2023judge} with explicit validation.
